%!TEX root = ../template.tex
%%%%%%%%%%%%%%%%%%%%%%%%%%%%%%%%%%%%%%%%%%%%%%%%%%%%%%%%%%%%%%%%%%%%
%% D1_abstract-pt.tex
%% NOVA thesis document file
%%
%% Abstract in Portuguese
%%%%%%%%%%%%%%%%%%%%%%%%%%%%%%%%%%%%%%%%%%%%%%%%%%%%%%%%%%%%%%%%%%%%

\typeout{NT FILE D1_abstract-pt.tex}%

A disseminação eficiente de dados em redes veiculares constitui um desafio significativo nos Sistemas de Transporte Inteligentes, devido à elevada mobilidade dos veículos e às rápidas mudanças de topologia. O paradigma \gls{NDN} apresenta-se como uma alternativa promissora às arquiteturas tradicionais baseadas em IP, ao centrar a comunicação nos próprios dados. Contudo, mecanismos de sincronização existentes, como o \gls{SVSync}, foram inicialmente concebidos para cenários estáveis e não tiram partido da dinâmica inerente às \glspl{VANETs}.

Neste trabalho, propõe-se um modelo \textit{hierárquico} de sincronização baseado no \gls{SVSync}, adaptado às características de redes veiculares \gls{V-NDN}. A arquitetura introduz três fases de sincronização suportadas por pivôs distribuídos, com o objetivo de reduzir a sobrecarga de mensagens, melhorar a escalabilidade e acelerar a convergência dos estados entre veículos. A solução foi implementada e avaliada no simulador \textit{ndnSIM}, onde foram simulados diferentes densidades de nós e analisadas métricas como latência, tráfego de sincronização, perdas de pacotes e escalabilidade.

Os resultados demonstram que o modelo hierárquico proposto melhora significativamente o desempenho da sincronização quando comparado com abordagens não hierárquicas, reduzindo o número de mensagens trocadas e aumentando a eficiência global. Estas conclusões contribuem para o desenvolvimento de soluções mais eficazes de disseminação e coordenação de dados em redes veiculares baseadas em \gls{NDN}.

% Palavras-chave do resumo em Português
\keywords{
  \gls{NDN} \and
  \gls{SVSync} \and
  Sincronização \and
  Redes Veiculares \and
  \glspl{VANETs}
}

